\section{Metodología}
    Para el desarrollo del proyecto se ha seguido un modelo de desarrollo estructural en cascada. 
    Este enfoque divide el proceso de desarrollo en fases claramente diferenciadas que se suceden de forma secuencial, comenzando cada fase una vez finalizada la anterior. 

    Las fases típicas de este modelo incluyen el análisis, el diseño, la implementación, las pruebas y el mantenimiento. 
    No obstante, al tratarse de un proyecto académico, la fase de mantenimiento no es necesaria y se sustituye por la documentación final de la memoria.

    Se emplea esta metodología debido a que el objetivo final del proyecto está claramente definido desde el inicio, con unos requisitos bien establecidos y sin necesidad de iteraciones incrementales. 
    Por ello, la planificación del proyecto se ha realizado tomando como base el análisis inicial, permitiendo organizar de forma coherente y estructural el desarrollo del trabajo.

\section{Planificación temporal}
    La organización temporal del proyecto se ha realizado siguiendo la dedicación exigida por la guía docente de la asignatura \cite{guiadocente-tfg}.
    El proyecto se ha planificado para ser elaborado en un total de 300 horas, distribuidas en 15 semanas, dedicando aproximadamente 20 horas semanales al desarrollo del mismo. 
    Se ha repartido el trabajo en 6 fases, estructuradas temporalmente según la Tabla \ref{tab:3.1}.

    \begin{table}[H]
        \centering
        \begin{tabular}{|c c c >{\centering\arraybackslash}p{8cm}|}
        \hline
        \textbf{Fase} & \textbf{Semanas} & \textbf{Horas} & \textbf{Entregables} \\ \hline
        Análisis & 1 & 10 & Estudio del alcance de la aplicación, definición de requisitos y casos de uso\\
        Planificación & 1 & 10 & Planificación temporal y seguimiento del proyecto\\
        Diseño & 2 & 20 & Decisiones técnicas, diagramas UML, diseño de la base de datos y mockups\\
        Implementación & 3-8 & 120 & Sistema funcional\\
        Pruebas & 8-9 & 40 & Validación del sistema\\
        Documentación & 10-15 & 100 & Memoria final y presentación\\
        \hline
        \end{tabular}
        \caption{Planificación temporal inicial del proyecto.}
        \label{tab:3.1}
    \end{table}

\section{Estimación de costes}
    La estimación de costes se ha realizado considerando el tiempo de dedicación al proyecto y el coste asociado a la hora de trabajo.
    El coste de un desarrollador junior se sitúa aproximadamente en 15€/h, lo que supone un coste total estimado de 4500€ para las 300 horas de trabajo planificadas. 

    Durante la fase de desarrollo no se prevén costes adicionales, ya que se utilizarán herramientas gratuitas y servicios que ofrecen planes con créditos gratuitos suficientes para cubrir las necesidades del proyecto.

\section{Análisis de riesgos}
    En esta sección se muestran los principales riesgos identificados durante la planificación del proyecto, junto con su probabilidad de ocurrencia, su impacto y las posibles medidas para mitigarlos, Talba \ref{tab:3.3}. 

    \begin{table}[H]
        \centering
        \begin{tabular}{| p{4.5cm} | c | c | p{7cm} |}
            \hline
            \textbf{Riesgo} & \textbf{Probabilidad} & \textbf{Impacto} & \textbf{Salvaguarda} \\ \hline
            Créditos insuficientes en servicios en la nube & Media & Alto & Monitorización del gasto de créditos y evitar el uso innecesario de recursos. \\ \hline
            Incremento inesperado del coste de las tecnologías utilizadas & Baja & Medio & Evaluación continua de los costes asociados y búsqueda de alternativas gratuitas o de bajo coste. \\ \hline
            Indisponibilidad temporal de los servicios externos & Baja & Alto & Uso de servicios consolidados con buen soporte. \\ \hline
            Aparición de imprevistos que afecten al tiempo de desarrollo & Alta & Medio & Planificación de tiempo adicional para imprevistos y margen con la fecha límite de entrega. \\ \hline
            Complejidad técnica en la integración de los distintos componentes del sistema & Media & Medio & Formación y estudio previo de las tecnologías utilizadas. \\ \hline
        \end{tabular}
        \caption{Análisis de riesgos del proyecto.}
        \label{tab:3.2}
    \end{table}

\section{Seguimiento del proyecto}
    A lo largo del desarrollo del proyecto, se ha llevado un seguimiento detallado para controlar el gasto temporal real frente al planificado. 
    En la Tabla \ref{tab:3.3} se muestra la comparativa entre la planificación temporal inicial y el progreso real del proyecto, permitiendo identificar desajustes y calcular el coste real del proyecto.
    \begin{table}[H]
        \centering
        \begin{tabular}{|c c c | c|}
            \hline
            \textbf{Fase} & \textbf{Horas Programadas} & \textbf{Horas Reales} \\ \hline
            Análisis & 10 & 15 \\
            Planificación & 10 & 8 \\
            Diseño & 20 & 25 \\
            Implementación & 120 & 180 \\
            Pruebas & 40 & 35 \\
            Documentación & 100 & 90\\
            \hline
        \end{tabular}
        \caption{Seguimiento del proyecto.}
        \label{tab:3.3}
    \end{table}
    El gasto temporal real del proyecto asciende a 353 horas por imprevistos en la fase de implementación. 
    
    Inicialmente, estaba previsto el uso de la red telefónica para la comunicación entre los usuarios y el sistema, lo que no aumentaba el gasto temporal ya que venía integrado con los servicios de la IA.
    Sin embargo, debido al elevado coste de esta solución en desarrollo, se optó por una solución basada en llamadas VoIP, que resulta completamente gratuita, pero implica un desarrollo adicional para integrar esta funcionalidad en el sistema.

    En cuanto al coste real del proyecto, el gasto ascendería a 5295€, superando el presupuesto inicial en 795€. 
    El gasto de la red telefónica hubiera supuesto un coste adicional de 1,50€ de establecimiento de llamada más 1€ adicional por cada minuto. 
    Además, el coste de alquiler de una línea telefónica virtual hubiera implicado un gasto mensual de entre 5-10€/mes. 
    Estos gastos hubieran complicado el desarrollo del proyecto, limitando el número de llamadas que se podrían realizar durante la fase de implementación y pruebas, lo que hubiera afectado negativamente a la calidad del sistema.